\clearpage
\chapter*{Đính chính cho Bài trình bày số 195}
\addcontentsline{toc}{chapter}{Đính chính cho Bài trình bày số 195}
\setcounter{footnote}{0}
\sid{e195}
\markboth{Đính chính 195}{Đính chính 195}
\src{300}

\noindent\textbf{4. TDTE II: Định lý tồn tại trong lý thuyết hình thức về
môđuli.}

\begin{center}
[Hội thảo Bourbaki, tập 12, 1959/60, số 195, 22 tr.]
\end{center}

\noindent\textbf{Trang 8.} --- Công thức viết trong Mệnh đề 5.1 chỉ đúng khi
$\Lambda$ là một trường; trong trường hợp tổng quát, phải thay
$\mathfrak m_\xi/\mathfrak m_\xi^2$ bằng thương của không gian này theo ảnh
của $\mathfrak n_\xi/\mathfrak n_\xi^2$, trong đó $\mathfrak n$ là iđêan cực
đại của $\Lambda$. Hơn nữa, định nghĩa đã cho về việc $O$ trơn trên
$\Lambda$ chỉ đúng khi mở rộng trường thặng dư $k'/k$ khả phân. Trong trường
hợp tổng quát, xem SGA III, 1.1.

\noindent\textbf{Trang 14, nhận xét.} --- Các bài toán nêu ở đây được giải
quyết hoàn toàn trong trường hợp xạ ảnh bằng các ``lược đồ Hilbert'' (xem
TDTE IV). Trái lại, những ví dụ của Nagata và Hironaka cho thấy các hàm tử
được xét không nhất thiết biểu diễn được nếu không áp đặt giả thiết xạ ảnh,
ngay cả khi chỉ phân loại các đa tạp con chiều không của một đa tạp đầy đủ,
không kỳ dị, chiều 3; điều này liên hệ với sự kiện rằng bình phương đối xứng
của một đa tạp như thế có thể không tồn tại.

\src{301}
\noindent\textbf{Trang 15, số 3.} --- Để có một nghiên cứu đầy đủ hơn, xem
TDTE V.

\noindent\textbf{Trang 16, Mệnh đề 3.1.} --- Thay vì ``nếu $f$ riêng'',
hãy đọc ``nếu $f$ riêng và khả tách''.

\noindent\textbf{Trang 18 và 19, các nhận xét, số 2.} --- Gần đây tôi đã
chứng minh rằng lược đồ hình thức môđuli của một đa tạp abel trên một trường
thực sự trơn trên vành Witt; nói cách khác, mọi lược đồ abel trên một vành
Artin địa phương vốn là thương của một vành khác đều thu được, bằng phép rút
gọn, từ một lược đồ abel trên vành sau. Chứng minh chỉ sử dụng các tính chất
biến thiên của lớp cản trở đối với phép nâng đã được đưa vào trong Bài trình
bày số 182, trang 12. Cũng xin nhắc rằng các lược đồ môđuli của các đường cong
giống $g$ và của các lược đồ abel phân cực đã được Mumford xây dựng (xem
SHMT).

\noindent\textbf{Trang 19, đoạn 5.} --- Phải giả sử tiết diện được xét để định
nghĩa $U$ không là ước của không chỉ trên $\mathfrak X$ mà còn trên mọi $X_n$.

\noindent\textbf{Trang 20, dòng 12 và 13.} --- Thay vì
\[
  G_x=G\bigl(\operatorname{Spec}(\widehat{\mathcal O}_{\mathfrak X,x})\bigr),
\]
hãy đọc
\[
  G_x=G\bigl(\operatorname{Spec}(\mathcal O_{\mathfrak X,x})\bigr)
  \subset
  G\bigl(\operatorname{Spec}(\widehat{\mathcal O}_{\mathfrak X,x})\bigr),
\]
và, thay vì ``$\mathcal O'_{0x}$'', hãy đọc
``$\widehat{\mathcal O}'_{0x}$''.

\par
